\section{Objetivo general}
Implementar un sistema de secuenciación inmersiva y sincronizada mediante una Arquitectura Dirigida por Eventos (EDA) en un videoclip musical interactivo de artistas guatemaltecos en un entorno de realidad virtual (VR) para analizar la percepción de presencia de usuarios mediante el \emph{I-Group Presence Questionnaire} (IPQ).

\section{Objetivos específicos}
\begin{itemize}
\item Diseñar la arquitectura del sistema de secuenciación bajo un enfoque EDA, definiendo productores, consumidores, eventos y gestor de eventos para coordinar elementos visuales, auditivos e interactivos en un entorno de VR.
\item Desarrollar un sistema de ejecución de eventos temporalizados que sincronice la reproducción del videoclip con la activación de escenas, iluminación, elementos visuales, audio e interacciones en momentos específicos de la canción.
\item Evaluar la percepción de presencia de los usuarios en la experiencia desarrollada mediante IPQ, considerando las dimensiones de presencia espacial, involucramiento y realismo

\end{itemize}